%=============================================================================
%  chapter7.tex -- "What It Means"
%
%  BEAT: hand over something usable.
%
%  Entry state: the reader has the result (ch5), its mechanism, its bound, and
%  every objection that could be computed inside a 0-D atomic model (ch6).
%  Chapter question: what should a diagnostician do differently, and what may
%  an atomic physicist claim is new?
%  Contribution: a procedure at an arbitrary operating point; the smallest
%  defensible novelty statement; the four open questions worth the cost.
%  Exit: none. This is the last chapter.
%
%  RULE observed here: nothing is re-derived. Every number names the chapter,
%  script or artifact it came from.
%=============================================================================

\chapter{What It Means}
\label{ch:conclusions}

The physicist of Chapter~\ref{ch:intro} is still at her console. Nothing in this
thesis takes her spectrometer away, alters the two wavelengths she measures, or
questions the ratio she forms from them. What it changes is the moment at which
she is entitled to look that ratio up. This chapter states that condition, hands
over the two numbers that convert it into an error estimate for her own
transient rather than for the one modelled here, and says how much of the
reasoning behind it belongs to this work and how much to a literature that is
sixty years old.

%=============================================================================
\section{The suspect was innocent}
\label{sec:what_found}
% QUESTION: what did the thesis set out to test, and what did it find instead?
%=============================================================================

Chapter~\ref{ch:intro} asked how large the error in a $(\Te,\nel)$ inversion of
the $H_\alpha/H_\beta$ ratio becomes when the plasma is disturbed faster than
its ionisation balance can relax, what physical quantity sets that error, and
where in the operating range it is largest. It also named a suspect. The
quasi-steady-state (QSS) approximation eliminates several dozen dynamical
variables by asserting that the excited levels equilibrate faster than
conditions change, and the expectation was that some of them would fail to keep
up.

They keep up. Integrating the full 43-state system through a temperature step
and evaluating both error measures along the same trajectory gives
$\epsQSS = 8.66\times10^{-6}$ at the benchmark point and
$\epsQSS = 6.73\times10^{-9}$ at $[0,4]$, the grid point carrying the largest
diagnostic error anywhere (Section~\ref{sec:qss_is_exact}). The closure is at
its most accurate exactly where the tabulated answer is at its worst. There is
no intermediate eigenvalue for an excited population to lag on: the gap between
the slow mode and the fast band is $3.7\times10^{7}$ at that point, so a level
that responds in nanoseconds cannot lag a reservoir that moves in microseconds.

What fails is the other assumption, the one Section~\ref{sec:the_assumption}
separated out and then left standing. A table with two axes has to dispose of a
third argument, the neutral-to-ion ratio $u = \ngs/\nion$, and the rule it uses
to dispose of it is equilibrium ionisation balance. That rule asserts that the
plasma has finished changing. For the whole of $\tauqss$ after a disturbance it
is false, and the excited manifold, having re-equilibrated in nanoseconds onto a
reservoir that has not moved, reports the partial-equilibrium state faithfully.
The ratio is measured correctly and the table then assigns it to the wrong
plasma.

If the fast states were the problem, the repair would be a time-dependent solve
of the excited manifold, and Section~\ref{sec:qss_is_exact} shows that such a
solve returns an answer already correct to eight decimal places. The repair
belongs instead at the closure applied to $u$, which is a statement about the
plasma's recent history and not about its atoms. Calling the failure a
breakdown of the quasi-steady-state approximation would point the effort at the
one part of the calculation that is working.

%=============================================================================
\section{What to do at the console}
\label{sec:criterion}
% QUESTION: given a measured ratio and a plasma that has recently been
% disturbed, what should a diagnostician compute before trusting the table?
%=============================================================================

The first thing to say is what the error is not. It is not a bias that can be
determined once and subtracted. Its size depends on the transient the plasma has
just undergone, and a single line ratio contains no record of that transient:
the same measured number is returned by an equilibrated plasma and by a
displaced one. No amount of care with the spectrometer separates them, and no
correction factor tabulated on $(\Te,\nel)$ can carry the information, because
the information is not a function of $(\Te,\nel)$.

What can be tabulated is the plasma's \emph{response} to being displaced, and
that is a property of the operating point alone. Section~\ref{sec:exact_decomposition}
factorises the plateau error exactly into a sensitivity and a displacement, and
Equation~\eqref{eq:gain_def} splits the displacement again into a gain and
the size of the transient, leaving
\begin{equation}
  \epsplat \;=\; \Big|\exp\!\big(\overline{S}\,G\,\Delta\ln\Te\big)-1\Big| ,
  \qquad
  \lim_{\Delta\ln\Te\to0}\frac{\mathrm{d}\,\epsplat}{\mathrm{d}\ln\Te}
  \;=\; \big|\overline{S}\,G\big| .
  \label{eq:recipe}
\end{equation}

\noindent This repeats Eq.~\eqref{eq:eps_of_step}, and it is repeated because
its function has changed. In Chapter~\ref{ch:results} it was the end of a
derivation; here it is a procedure, and the reader supplies the third factor. It
inherits that chapter's assumption that the ion population is a fixed reservoir
and that $\ngs$ is frozen for the duration of the event, which removes the
freedom for ionisation balance to respond and is the condition
Section~\ref{sec:transport} cannot test.

Both structural factors are dimensionless. The mean sensitivity $\overline{S}$
is the average of $\ffrac{3}-\ffrac{4}$ over the interval the reservoir
traverses, and it says how strongly the observed ratio answers to the
ground-state reservoir. The reservoir gain $G = \partial\ln u^{\mathrm{CRE}}/\partial\ln\Te$
says how far equilibrium ionisation balance would move that reservoir for a
given fractional change in temperature; it is negative everywhere on the grid,
because heating ionises the ground state away.

\paragraph{The procedure.} For a CR matrix of one's own, at one's own operating
point:

\begin{enumerate}
\item Solve the excited manifold twice, once with the ion reservoir alone and
once with the ground state alone, to obtain the population coefficients $a_p$
and $c_p$ of Eq.~\eqref{eq:two_supplies_intro} for the two emitting shells.
These are the same two solves any steady-state CR table already performs. Form
the switching points $x_m = \ln(c_m/a_m)$ and their separation
$\Delta = x_4 - x_3$; the two ground-fed fractions are then the single logistic
of Eq.~\eqref{eq:f_logistic}, displaced by $\Delta$.
\item Differentiate the equilibrium reservoir with respect to temperature at
fixed density to get $G$.
\item Insert the transient one actually cares about as $\Delta\ln\Te$. The
reservoir excursion is $\Delta\ln u = G\,\Delta\ln\Te$, and that fixes the
interval over which $\overline{S}$ is the mean of $\ffrac{3}-\ffrac{4}$.
\end{enumerate}

Steps 1 and 2 belong to the operating point and are done once; step 3 belongs to
the event. The sensitivity straddles the two, because the curves are the
operating point's and the interval is the event's, and in the small-step limit
it collapses onto the operating point as well: $\overline{S}$ tends to
$\ffrac{3}-\ffrac{4}$ evaluated at the undisturbed reservoir, which is why
$|\overline{S}G|$ can be tabulated on $(\Te,\nel)$ alone. That division is the
whole reason for reporting structure instead of a percentage. Measured across
the grid, $|G|$ moves by at most \SI{6.3}{\percent}
between step sizes of one, two and four grid intervals, over the $736$
point-and-direction triples for which all three step sizes were computed, while
$\epsplat$ at the same points changes by up to a factor $8.44$. A percentage
quoted without its step is not a property of the plasma; $G$ very nearly is.
Both figures are direct reductions of
\texttt{validation/reservoir\_gain/reservoir\_gain.csv} ($2288$ rows, stamped
10 September 2026). \emph{[UNVERIFIED: no script in the repository performs
these two aggregations. What would close this: a reduction script that writes
the spread of $|G|$ and the range of $\epsplat$ over $k$ to \texttt{validation/}
alongside the table it summarises.]}

\paragraph{The scale, before anyone computes anything.} Over the $982$ heating
steps of that file that have a resolvable plateau window, $|\overline{S}|$ runs
from $0.0649$ to $0.4822$; over all $2288$ rows $|G|$ runs from $2.640$ to
$14.519$. Restricted to the range this thesis defends, $\Te\ge\SI{2}{\electronvolt}$,
those become $0.0649$ to $0.4552$ and $2.640$ to $7.822$, and the product
$|\overline{S}G|$ runs from $0.362$ to $3.535$. The upper end is the useful
number to remember. In the worst conditions defended here, a \SI{1}{\percent}
temperature excursion costs the Balmer inversion about \SI{3.5}{\percent} in the
line ratio; at the benchmark point it costs about \SI{1.3}{\percent}.

\subsection{Four things that follow, and one that does not}

\paragraph{The response has a ceiling, and it is below unity.} Because both
ground-fed fractions are the same unit-width logistic curve displaced by
$\Delta$, their difference cannot exceed $\tanh(|\Delta|/4)$, and $\overline{S}$
inherits that bound (Eq.~\eqref{eq:Sbar_bound}). The $\tanh$ comes from the
width of the logistic, and the width is unity because the excited populations
are exactly affine in $\ngs$, which is a property of the linear system rather
than of any rate coefficient. Every property of the 42-state network enters only
through $\Delta$. Since $|\mathrm{d}\ln\Robs/\mathrm{d}\ln\bdep| < 1$ everywhere,
integrating over any excursion gives $|\Delta\ln\Robs| \le |\Delta\ln\bdep|$, so a
reservoir left stale by a factor of two cannot put the line ratio out by more
than a factor of two. A stale reservoir is attenuated in logarithmic space,
never amplified. That protection holds for every shell pair, every
$(\Te,\nel)$ and every atomic dataset, and it is thinnest where the error is
largest: at $[23,3]$, one column into the crest, the measured sensitivity
already sits at \SI{97}{\percent} of the ceiling (Section~\ref{sec:bound}), so
no refinement of the rate coefficients can move it far in either direction.

\paragraph{The clock is $\tauqss$, and it is far too slow to average away.} The
temptation is to assume that an error introduced by a fast disturbance decays
with the fast timescale. It does not. The error decays as a single exponential
on $\tauqss$, with a half-life of \SI{12.95}{\micro\second} at the benchmark and
\SI{1.10}{\milli\second} at $[15,3]$ (Section~\ref{sec:persistence}), against
the \SI{2.28}{\nano\second} on which the excited manifold responds. A
\SI{100}{\micro\second} exposure does not integrate the error away.
Counted on the time-averaged bound of Eq.~\eqref{eq:elm_bound}, which direct
43-state propagation places within \SI{0.7}{\percent} of the true average, $45$
of $448$ point-and-direction pairs above \SI{2}{\electronvolt} exceed
\SI{10}{\percent}, the worst being \SI{17.5}{\percent} at
$\Te = \SI{2.02}{\electronvolt}$ and $\nel = \SI{1.93e13}{\per\cubic\centi\metre}$.
At densities where the Balmer inversion is ordinarily applied,
$\nel\ge\SI{e14}{\per\cubic\centi\metre}$, none of $108$ pairs exceed
\SI{10}{\percent} and the worst is \SI{7.2}{\percent}. The conditions to
distrust are the tenuous and cool ones, and the error grows toward low density
and low temperature together.

\paragraph{Timescale separation gives no warning.} A modeller deciding whether
to trust a QSS-based inversion reaches for $\Msep = \tauqss/\taurel$, which is
the ratio that licenses eliminating the fast variables in the first place. It
carries no information about the diagnostic error once the operating point is
controlled for (Section~\ref{sec:M_does_not_predict}). The two maxima are
$52\times$ apart in density; at the point of largest $\Msep$ the plateau error
is \SI{12.0}{\percent}, the 75th percentile of its distribution and a factor
$3.2$ below the worst. The reason is visible in Eq.~\eqref{eq:recipe}. $\Msep$
compares two rates, and $\epsplat$ measures a distance, the product of a
sensitivity with a displacement, neither of which appears in $\Msep$. A large
$\Msep$ certifies that the excited states have equilibrated with the reservoir.
It says nothing about whether the reservoir is in the right place.

\paragraph{A second failure, unrelated to any of this.} The bound of
Eq.~\eqref{eq:Sbar_bound} constrains the error in the ratio, and a
diagnostician does not suffer an error in the ratio; she inverts it. The
inversion divides by $\partial\ln\Robs/\partial\ln\Te$, and that derivative
passes through zero. Scanning temperature at
$\nel = \SI{5.18e13}{\per\cubic\centi\metre}$, the fractional distance
$\epsstep$ between the two tabulated answers either side of the step falls to
$4.9\times10^{-5}$ near \SI{6.9}{\electronvolt}
(Section~\ref{sec:inversion_jacobian}), which means the tabulated ratio is
stationary in temperature there. The $n=3/n=4$ ratio is locally non-invertible
for temperature at that point, so a bounded error in the observable produces an
unbounded error in the inferred temperature. This has nothing to do with
transients and would afflict a perfectly equilibrated plasma. For a practitioner
it is the more dangerous of the two failures, and it was found only because
$\epsstep$ was computed and kept rather than discarded. Its locus across the
rest of the grid has not been mapped, and Section~\ref{sec:inversion_jacobian}
records what would close that.

\paragraph{What does not follow.} None of the above licenses a claim about a
particular machine. Every number in this section is conditional on the
ground-state neutral density being frozen for the duration of the event, which
is the conditional Section~\ref{sec:transport} states and cannot test from
inside a zero-dimensional model.

%=============================================================================
\section{What is new, at the smallest size that is defensible}
\label{sec:novelty}
% QUESTION: which parts of the argument are this work's, and which were
% established decades ago?
%=============================================================================

Section~\ref{sec:prior_work} set out what was already known before this work
began. What follows is the narrower question of what may now be claimed that
could not be claimed then, item by item, and it is deliberately smaller than the
argument might suggest.

\paragraph{The two-channel decomposition is not new and is not claimed.} That
every excited population is exactly affine in the ground-state and ion
densities, Eq.~\eqref{eq:two_supplies_intro}, is the population-coefficient
decomposition of Bates, Kingston and McWhirter, written in the Saha-normalised
form $b(p) = r_0(p) + r_1(p)\,b(1)$ by Fujimoto and McWhirter
\cite{Bates1962, FujimotoMcWhirter1990, Fujimoto2004}. Chapter~\ref{ch:theory}
says so where it uses it. Everything structural in this thesis is a consequence
of that 1962 identity.

\paragraph{The logistic form is a rewriting, and should be called one.} Writing
the ground-fed fraction in the logarithmic reservoir variable turns
Eq.~\eqref{eq:ground_fed_fraction} into the single logistic of
Eq.~\eqref{eq:f_logistic}, so the two shells' fractions are the same curve
displaced. This is one line of algebra applied to a published formula, and the
same function appears in chemical kinetics as the Michaelis--Menten saturation
fraction, where a substrate concentration plays the part of the reservoir. The
defensible statement is that it has not previously been stated in this form in
the collisional--radiative literature. It is not a new function and it is not a
new result about hydrogen.

\paragraph{The bound is elementary; its application is the claim.} That $\max|\ffrac{3}-\ffrac{4}| = \tanh(|\Delta|/4) < 1$ follows in a
few steps from two displaced logistics, and no prior statement of it was found.
That absence is weak evidence. The literature search behind it reached Crossref
and OpenAlex, which index journal metadata well and equations buried inside
textbook chapters poorly, and this identity is exactly the kind of relation that
would sit unlabelled in a monograph. What this work claims is the
\emph{application}: that the inequality bounds the diagnostic error of
a hydrogen line ratio with respect to its ground-state reservoir, uniformly over
every atomic dataset, and that the bound is near-tight where the error is
largest. The novelty of the inequality itself is left open.
\emph{[UNVERIFIED: two bodies of work must be read before the novelty question
can be settled, and neither has been read against it. Fujimoto's \emph{Plasma
Spectroscopy} Chapter~4 \cite{Fujimoto2004}, which develops the population
coefficients in the normalised form this work uses; and Fujimoto's series on
the kinetics of ionisation and recombination in the Journal of the Physical
Society of Japan \cite{Fujimoto1979a, Fujimoto1979b, Fujimoto1980a,
Fujimoto1980b, Fujimoto1985}, which runs to five parts rather than the four it
is usually cited as, and which is where a closed-form statement about the ratio
of two population coefficients would most plausibly already appear. The
bibliographic records for all five parts were verified against Crossref and
OpenAlex on 10 September 2026; the papers themselves have not been obtained.]}

\paragraph{Greenland stated the negative result in general form, twenty-five
years ago.} That the validity criteria for a reduced CR model are conditions on
eigenvectors rather than eigenvalues, and are not related to the equilibrium
timescales, is Greenland's, in two papers of the same year and separately
attributable \cite{Greenland2001a, Greenland2001b}. This work does not improve
on either statement. It quantifies them for one diagnostic and maps where the
criterion misleads: the $52\times$ separation of the two maxima, the sign
reversal of the association once the operating point is controlled for, and the
census of Section~\ref{sec:M_does_not_predict}. Turning a general criterion into
a number is a smaller contribution than stating the criterion, and it is the one
available here.

\paragraph{Sawada and Fujimoto asked the adjacent question and answered it
first.} Their paper is titled, in part, for the validity range of the
quasi-steady-state solution of the coupled rate equations
\cite{SawadaFujimoto1994}, and it treats the same system: atomic hydrogen, a CR
model, an abrupt change of conditions. What they established is which level
controls the approach to steady state. In an ionising plasma the levels below
Griem's boundary relax at their own natural lifetimes while those above relax at
the rate of the boundary level, and in a recombining plasma the roles invert.
Their answer is a clock internal to the excited manifold. The question here is
what a diagnostic inversion costs while the reservoir underneath that manifold
has not settled, which is a different quantity governed by a clock four orders
of magnitude slower, and which can be large at the very moment the excited-level
solution is exact. The two results are compatible and neither contains the
other.

\paragraph{What is claimable.} The cost of the ionisation-balance
closure hidden inside a two-parameter Balmer table is now a computable number at
an arbitrary operating point, expressed through two structural coefficients that
are properties of the plasma rather than of the disturbance, together with a
ceiling on it that no atomic dataset can raise and a map of where in
$(\Te,\nel)$ the coefficients are large. That is the contribution, and it is
smaller than the question Chapter~\ref{ch:intro} asked.

%=============================================================================
\section{What this work cannot claim}
\label{sec:implications}
% QUESTION: what would an examiner be entitled to strike out?
%=============================================================================

These are carried forward from Chapter~\ref{ch:limits} without re-argument, and
they bound every statement in Section~\ref{sec:criterion}.

\paragraph{Everything is conditional on a frozen reservoir.} If neutral
transport and recycling renew $\ngs$ faster than the local atomic physics does,
the ground state is not stale and the error mechanism does not operate. The
model measures the atomic clock and assumes the transport one does not interfere
(Section~\ref{sec:transport}).

\paragraph{It is not a divertor result below two electron-volts, and the
restriction is not a hedge.} Holding $\nel$ fixed while the ground-state
population changes is self-consistent only where the plasma is largely ionised.
Above \SI{2}{\electronvolt} conservation of nuclei across one temperature step
demands $|\delta\nel/\nel| < 10^{-3}$. Below it, $68$ of $392$ one-step
operators require $|\delta\nel/\nel| > \SI{10}{\percent}$ and $36$ require more
than \SI{100}{\percent} (Section~\ref{sec:quasineutrality}).
$\Te\ge\SI{2}{\electronvolt}$ is the only self-consistent region of the grid,
not a cautious subset of it, and the cold-edge numbers are asymptotic values of
a model outside its own domain.

\paragraph{It is not detachment.} The density crest sits a factor $5.2$ below
the detached band, and at citable detached densities the ELM-averaged bound
never reaches \SI{10}{\percent} (Section~\ref{sec:not_detachment}). The crest is
a property of the atomic system and needs no divertor geometry to produce it.

\paragraph{There are no molecular channels, in a gas that is largely
molecular.} A third supply feeding $n=3$ preferentially would change
$\ffrac{3}-\ffrac{4}$. Since $\epsplat$ is linear in that difference, taking
molecular emissivity fractions of $0.60$ and $0.30$ at $n=3$ and $n=4$ reduces
the cold-corner plateau error by a factor of $3$ to $5$
(\texttt{findings\_10}~ADDENDUM~D.4). That bound applies to the cold corner,
which lies outside the defended range in any case.
\emph{[UNVERIFIED: Section~\ref{sec:molecules} states that no bound on the
molecular effect can be constructed from the data in this repository, and
ADDENDUM~D.4 constructs one from the two references cited in that same section.
The two statements are inconsistent. This chapter quotes the bound because it
exists; the reconciliation belongs in Chapter~\ref{ch:limits}.]}

\paragraph{The density maximum is a range, not a density.} It lies between
$7\times10^{12}$ and \SI{5e13}{\per\cubic\centi\metre}, resolved to about a
factor of $3$, because every density row sits above \SI{90}{\percent} of its own
maximum across roughly a decade (Section~\ref{sec:results_summary}). Quoting a single density for it overstates what
the grid resolves. It is also the $(3,4)$ shell pair's maximum and not the
Balmer series' (Section~\ref{sec:ridge_conditional}).

%=============================================================================
\section{What would settle what is open}
\label{sec:further_work}
% QUESTION: which unfinished calculations would change what may be claimed?
%=============================================================================

Four calculations would change the standing of a result in this thesis. The
first two cost days. The last two are different projects.

\paragraph{A joint $(\Te,\nel)$ step map.} The maps in
Chapter~\ref{ch:results} step temperature at fixed density, and an edge-localised
mode (ELM) raises both. Joint steps give a comparable plateau error but a
time-averaged bound that can fall by a factor of $4$, because the higher density
shortens $\tauqss$; at the benchmark the bound falls from $0.0117$ to $0.0030$
(Section~\ref{sec:persistence}). The census of Chapter~\ref{ch:results} is a fixed-density
census, and an examiner who knows ELM physics will ask for the other one. It
costs the same as the map that already exists.

\paragraph{Reading Fujimoto's Table~4.1(b) from the book.} Chapter~\ref{ch:limits}
publishes an unexplained factor-$8$ disagreement with the tabulated ground-fed
population coefficient at $p=3$, which is one of the two coefficients the whole
mechanism runs on (Section~\ref{sec:r1_deficit}). One sentence from the
monograph settles whether that table's level structure is bundled in $n$ with
statistical $\ell$ populations or resolved in $\ell$, and therefore whether the
comparison tests this model or tests the $\ell$-closure inside it. Until that is
read, the benchmark adjudicates nothing.

\paragraph{A two-region or transport-coupled calculation.} The single conditional
on which all of Section~\ref{sec:criterion} rests is that $\ngs$ is frozen for
the duration of the event. Deciding it requires $\ngs$ to be supplied by a
recycling influx from an edge transport calculation rather than solved for from
$\Lmat\nvec + \Svec\nion = 0$, which is a different class of calculation and
not a refinement of this one. It is the item that would convert the result from
a statement about a class of CR model into a statement about a machine.

\paragraph{Molecular channels, which may change the form and not just the
coefficients.}
Adding molecules as states inside $\Lmat$, fed from the same two reservoirs,
leaves the structural results untouched, and Section~\ref{sec:molecules} says
so. Adding molecular hydrogen as a \emph{third independent reservoir} does not.
The population balance then reads
$n_p = a_p\ngs + c_p\nion + m_p n_{\mathrm{H_2}}$, with $n_{\mathrm{H_2}}$ the
molecular density in \si{\per\cubic\centi\metre} and $m_p$ the dimensionless
number of atoms in level $p$ supported per molecule. The observable then
depends on two independent reservoir ratios rather than one, and the ground-fed fraction is no longer a single logistic in a single
variable. The bound of Eq.~\eqref{eq:Sbar_bound} is a statement about a
one-dimensional family of curves and there is no reason to expect it to survive
in that form. Which of the two descriptions a real detaching divertor requires
is a physics question this work does not answer.

\medskip

The physicist at the console keeps her spectrometer, her ratio and her table.
What she gains is a test she can apply before she trusts the answer: evaluate
$\overline{S}$ and $G$ at the operating point she is looking at, multiply by the
fractional temperature excursion her plasma has just undergone, and compare the
result with the precision she needs. If the plasma has been quiet for several
times $\tauqss$, the table is telling her about the plasma. If it has not, the
table is telling her about a plasma that no longer exists, and
Eq.~\eqref{eq:recipe} says by how much.
